\section{Legal proceedings beyond the booked liability}
\label{sec:apple-legal}

Knowing the companies within a group does not identify all the events that
could change their earnings. A legal proceeding can require a payment, but it
can also change how a company sells its products or earns licensing revenue.
Apple's reported App Store injunction illustrates the second channel: it
concerns developers' external purchase links on the U.S. storefront, rather
than a disclosed amount of damages. A forecast that represents legal exposure
only as a liability balance would leave that operating consequence outside
the model.

The annual Item 3 and interim Part II Item 1 each discuss Digital Markets Act
investigations, the Department of Justice lawsuit, Epic Games and other legal
proceedings. These are dated issuer reports, with different kinds of statements
inside each matter: allegations, reported decisions, pending motions,
management's position and possible consequences. Keeping those distinctions
allows a model to use what is reported without deciding the case itself
\citep[Item 3, printed p.~18]{apple2024}
\citep[Part II, Item 1, printed p.~18]{apple2025q1}.

\subsection{A loss assessment answers a narrower question than exposure}

Annual Note 12 and interim Note 9 give the same qualified assessment:
``In the opinion of management, there was not at least a reasonable possibility
the Company may have incurred a material loss, or a material loss greater than
a recorded accrual, concerning loss contingencies for asserted legal and other
claims'' \citep[Note 12, printed p.~46]{apple2024}
\citep[Note 9, printed p.~11]{apple2025q1}. The sentence identifies whose
assessment it is, concerns an incurred material loss, and allows for a recorded
accrual. It does not disclose a total legal accrual, allocate an accrual to one
of the named matters, or assert that all exposure is zero. Nor does it supply
a numerical probability or materiality threshold.

The annual MD\&A separately reports Apple's recognition policy: the company
records a liability when it is probable that a loss has been incurred and the
amount is reasonably estimable, a determination requiring significant judgment
\citep[Legal and Other Contingencies, printed p.~26]{apple2024}. That policy
explains why a pending proceeding and a recognised liability are different
facts. It is issuer-policy evidence; applying the operative US GAAP requirements
to a particular case still needs the relevant authority and case facts.
An allegation cannot supply those facts, and a possible change in future
sales is not itself the incurred loss described by the policy.

Settlement adds a further distinction. Annual Item 3 says certain matters
settled during the fourth quarter of fiscal 2024 without an individually or
aggregate material impact on financial condition or operating results. The
interim passage makes that statement about the first quarter of fiscal 2025.
These are two period-specific assertions, not the same settlement observation
copied forward. Neither passage identifies the matters, settlement amounts,
payment dates or earlier accruals. Both retain the warning that adverse
resolutions above management's expectations could materially affect the
reporting period. A payment schedule or an expense roll-forward therefore
remains unresolved even though the settlement statements themselves are known.

\subsection{The remedy, jurisdiction and procedural stage matter}

The Digital Markets Act discussion reports two investigations opened by the
European Commission on 25 March 2024. Article 5(4) concerns developers'
communication and promotion of offers and conclusion of contracts with end
users; Article 6(3) concerns default settings, app uninstallation and a browser
choice screen on iOS. The Commission's 24 June announcement is described as
preliminary findings alleging an Article 5(4) breach and the opening of a third
investigation into new contractual requirements for third-party developers
and app marketplaces. Apple states that it believes it complies and continues
to change its compliance plan in response to Commission engagement. The
preliminary allegation and Apple's position remain separate assertions.

Both filings say that, if the Commission makes a final violation determination,
it can issue a cease and desist order and may impose fines up to 10\% of
Apple's annual worldwide net sales. They also say an order would apply
immediately during an appeal unless a stay is granted. This preserves a
conditional consequence, its stated sales base and the stay qualification.
The percentage is not a fine already imposed, an expected loss or a booked
liability. Even a conditional fine calculation would need the applicable
operative law, relevant sales period and base, and assumptions about the
determination and remedy. None is supplied by multiplying a convenient annual
sales figure by the reported cap.

The Department of Justice discussion reports the civil antitrust lawsuit
filed on 21 March 2024 in the District of New Jersey by the DOJ and state and
district attorneys general. It alleges monopolization or attempted
monopolization in the markets for ``performance smartphones'' and
``smartphones.'' The DOJ seeks equitable relief; related state and federal
civil litigation seeks monetary damages and other nonmonetary relief. Apple
states that it has substantial defenses and intends to defend itself.
Treating all these proceedings as one quantified damages claim would erase
both the different requested relief and the absence of disclosed amounts.

Epic's case contains a reported court finding as well as unresolved procedure.
Apple says the Northern District of California court found certain App Store
Review Guidelines provisions to violate California's unfair competition law
and issued an injunction against prohibiting developers' external purchase
links. The stated scope is the U.S. storefront of the iOS and iPadOS App Store.
Apple reports implementing a compliance plan and filing its statement on
16 January 2024. Epic's motion disputing that plan and seeking enforcement is
reported as pending, opposed by Apple; Apple also reports filing a motion to
narrow or vacate the injunction on 30 September 2024. These passages do not
establish that the disputed plan was judicially accepted or quantify its sales
effect. An EU-wide scenario cannot inherit the U.S. injunction merely because
both matters concern the App Store.

The interim bodies repeat the substantive DMA, DOJ and Epic account, with
acronym or trademark presentation differences. Repetition is evidence of what
Apple reported in each filing, not proof that every intervening legal fact was
unchanged. The September 30 motion also falls after the annual September 28
period-end. Event date, reporting period and the date on which the report
became available must therefore remain distinct in a historical information
query.

\subsection{A partner's remedy can change Apple's revenue}

The risk discussion connects these matters to operations. Both filings report
changes to U.S. developer communications and EU changes to distribution,
payment processing, fee structures and developer tools. They also describe
licensing revenue from Google and other search providers, whose arrangements
are subject to investigations and proceedings
\citep[Risk Factors, printed pp.~14--15]{apple2024}
\citep[Part II, Item 1A, printed pp.~19--20]{apple2025q1}.
The interim adds that Google was found to have violated U.S. antitrust laws on
5 August 2024 and that the DOJ proposed remedies ranging from restrictions on
licensing agreements to prohibiting commercial search-distribution terms with
Apple for up to ten years. Apple says that, if implemented, those remedies
could materially affect its ability to earn that licensing revenue. The
finding is an earlier event newly described in this retained comparison;
the remedies are proposed, not reported here as imposed.

That channel differs from an Apple fine. The affected economic state would
include the search agreement, its revenue terms, duration, geography and
substitution possibilities. Likewise, an App Store remedy requires the
affected storefront, developer terms, fees, customer purchases and behavioural
response. The filings identify these channels without supplying the contract
economics needed to calculate them. They do not justify applying one loss
percentage to all Services revenue or treating every customer as affected.

A conditional calculation should consequently specify which matter and
remedy it assumes, when that remedy takes effect, which contracts it changes
and how the parties respond. A monetary-loss calculation additionally needs
recognition facts, any opening accrual and the settlement and payment schedule;
expense recognition and cash payment can occur in different periods. The
State Aid tax charge and escrow settlement, developed in
Section~\ref{sec:apple-tax-settlement}, retain their own tax records and must
not be counted again as a DMA loss. Until the missing
facts and applicable treatments are supplied, the model can retain the reported
proceedings and affected queries as unresolved, without assigning probabilities
or producing numerical journals. This preserves useful knowledge about future
exposure while keeping the distinction between a known assertion and a
supported financial consequence.
